\chapter{Units and Prefixes}
\label{app:units}
A quick reference for exam conditions. \Cref{ch:foundations} develops where the base
relationships come from and why the ampere, not the coulomb, is the SI base unit, so
the first table below is deliberately a restatement---a lookup table is worth having
even when the reasoning lives elsewhere. The \emph{third} table is not a restatement:
it collects the compound and referenced units the later chapters introduce one at a
time and never gather in one place, which is the thing hardest to look up mid-problem.

\section{Electrical Units}
\begin{tabularx}{\textwidth}{@{}L{0.22\textwidth}L{0.18\textwidth}Y@{}}
\toprule
Unit & Symbol & Equivalent relationship \\
\midrule
coulomb & \si{\coulomb} & \(\SI{1}{\coulomb}=\SI{1}{\ampere\second}\). \\
ampere & \si{\ampere} & \(\SI{1}{\ampere}=\SI{1}{\coulomb\per\second}\). \\
volt & \si{\volt} & \(\SI{1}{\volt}=\SI{1}{\joule\per\coulomb}\). \\
ohm & \si{\ohm} & \(\SI{1}{\ohm}=\SI{1}{\volt\per\ampere}\). \\
siemens & \si{\siemens} & \(\SI{1}{\siemens}=\SI{1}{\per\ohm}\). \\
farad & \si{\farad} & \(\SI{1}{\farad}=\SI{1}{\coulomb\per\volt}\). \\
henry & \si{\henry} & \(\SI{1}{\henry}=\SI{1}{\volt\second\per\ampere}\). \\
watt & \si{\watt} & \(\SI{1}{\watt}=\SI{1}{\joule\per\second}=\SI{1}{\volt\ampere}\). \\
joule & \si{\joule} & Energy. \\
hertz & \si{\hertz} & Cycles per second. \\
\bottomrule
\end{tabularx}

\section{Prefixes}
\begin{center}
\begin{tabular}{@{}lll@{}}
\toprule
Prefix & Symbol & Factor \\
\midrule
pico & p & \(10^{-12}\) \\
nano & n & \(10^{-9}\) \\
micro & \(\mu\) & \(10^{-6}\) \\
milli & m & \(10^{-3}\) \\
kilo & k & \(10^{3}\) \\
mega & M & \(10^{6}\) \\
giga & G & \(10^{9}\) \\
\bottomrule
\end{tabular}
\end{center}

\section{Compound and Referenced Units Used Later}
These appear scattered through Parts~III--V. The point of tabulating them is that
several are \emph{referenced} rather than absolute---a number in dBm means nothing
until you know the reference---and two of them are easy to confuse.

\begin{tabularx}{\textwidth}{@{}L{0.20\textwidth}L{0.20\textwidth}Y@{}}
\toprule
Unit & Where used & What it means \\
\midrule
VA (\si{\volt\ampere}) & \cref{ch:ac} & Apparent power \(|S|\); equals watts only when the load is purely resistive. \\
VAR & \cref{sec:instpower} & Reactive power \(Q\); energy shuttled in and out of storage each cycle, dissipating nothing. \\
dBW, dBm & \cref{sec:decibels} & Power referenced to \SI{1}{\watt} and to \SI{1}{\milli\watt}; the two differ by exactly \SI{30}{\decibel}. \\
dBm/Hz & \cref{sec:noisefloor} & Power \emph{spectral density}---power per unit bandwidth, hence the \SI{-174}{dBm} floor. \\
dBi, dBd & \cref{sec:antennalength} & Antenna gain referenced to an isotropic radiator or to a dipole; \(\mathrm{dBi}=\mathrm{dBd}+2.15\). \\
Np (neper) & \cref{sec:lineloss} & Natural-log attenuation, the units of \(\alpha\); \(1\,\mathrm{Np}=\SI{8.686}{\decibel}\), the \(20\log_{10}\ee\) of that section. \\
ppm & \cref{sec:crystal}, \cref{sec:instrumentlimits} & Parts per million, \(10^{-6}\) fractional; the natural unit for crystal pulling and timebase error. \\
\si{\kelvin\per\watt} & \cref{sec:thermalrunaway} & Thermal resistance \(\theta_{JA}\); temperature rise per watt dissipated. \\
\si{\siemens\per\meter} & \cref{sec:selfresonance} & Conductivity \(\sigma\); copper is \SI{5.8e7}{\siemens\per\meter}. \\
\si{\henry\per\meter} & \cref{sec:selfresonance} & Permeability \(\mu\); \(\mu_0=4\pi\times10^{-7}\). \\
S unit & \cref{sec:sunits} & Receiver signal-strength step of \SI{6}{\decibel}, i.e.\ a factor of four in power. \\
\bottomrule
\end{tabularx}

\begin{mistakebox}
Two pairs cause most of the trouble. \textbf{VA and VAR are not watts}: only watts
leave the circuit as heat or radiation, and a purely reactive load can show large VA
and VAR while consuming no average power (\cref{sec:instpower}). And \textbf{a bare
``dB'' is not a level}---it is a ratio, so ``\SI{10}{\decibel}'' answers ``how many
times?'' while ``\SI{10}{dBm}'' answers ``how much?''. Adding two bare decibel figures
is meaningful; adding two dBm figures is not.
\end{mistakebox}
